The Codynamic Book Machine is an intent-driven, schema-backed, multi-agent authoring system for producing long-form technical writing as a coordinated workflow rather than a sequence of isolated drafts. It begins with author intent and outline structure, normalizes that intent into a canonical work object, and then distributes section-level labor across specialized agents that draft, review, revise, and assemble content into publication artifacts.

The central claim of the system is that a book can be managed as a living, machine-readable work graph in which structure, dependencies, messages, and generated artifacts evolve together. Instead of treating outline, prose, citations, diagrams, and build outputs as separate concerns, the Codynamic Book Machine keeps them aligned through shared schema, durable task state, and agent coordination. This makes the book both the product and the operating context of the system.

In practice, the machine supports a pipeline that starts with conversational intake or imported outline material, converts that input into a canonical representation, decomposes the work into section-sized tasks, and routes those tasks through agents responsible for drafting, reference handling, diagram generation, design, and assembly. The result is a structured authoring environment in which changes can be traced, dependencies can be managed explicitly, and publication artifacts can be regenerated from the same underlying work object.

This paper presents the architecture, data model, and coordination workflow of the Codynamic Book Machine, and uses the system itself as a working example of codynamic design: a writing system whose content and control structure are developed together.
